\chapterbackmatter{Supplementary Exercises}

\begin{enumerate}
  \SupplementaryExercise{1}{Counting a General Real Superfield}
    {(Adapted from Philip C. Argyres, \emph{An Introduction to Global
    Supersymmetry}, Section 1.5.2)}
    {argyres-global-susy-2001-ch26}
  Starting from the finite Grassmann expansion
  \eqref{eq:26.2.10}, count the independent real bosonic and fermionic
  components of a real scalar superfield.  Explain why the equality holds
  off shell, before any field equation or gauge invariance is imposed.

  \SupplementaryExercise{2}{Closure on the Wess--Zumino Scalar}
    {(Adapted from Philip C. Argyres, \emph{An Introduction to Global
    Supersymmetry}, Sections 1.5.3--1.5.5)}
    {argyres-global-susy-2001-ch26}
  Use Eq.~\eqref{eq:26.1.21} to evaluate
  \([\delta_\alpha,\delta_\beta]A\).  Show explicitly that the auxiliary
  fields cancel and that the result is a translation.  State the
  corresponding result for \(B,\psi,F,G\).

  \SupplementaryExercise{3}{The Product Rule for Lowest Components}
    {(Adapted from Adel Bilal, \emph{Introduction to Supersymmetry},
    Sections 4.1--4.3)}
    {bilal-introduction-susy-2001-ch26}
  Let \(S=S_1S_2\) be the product of two bosonic general superfields.
  Derive \(C^S,\omega^S,M^S,N^S,V_\mu^S\) directly from
  Eq.~\eqref{eq:26.2.10}, including all fermion signs.

  \SupplementaryExercise{4}{Why the Covariant Superderivative Works}
    {(Adapted from Stephen P. Martin, \emph{A Supersymmetry Primer},
    version 7, Sections 4.1--4.3)}
    {martin-susy-primer-v7-ch26}
  Starting from Eqs.~\eqref{eq:26.2.2} and \eqref{eq:26.2.26}, verify
  \(\{\SuperD_\alpha,\SuperQ_\beta\}=0\) and
  \(\{\SuperD_\alpha,\bar{\SuperD}_\beta\}
  =-2\gamma^\mu_{\alpha\beta}\partial_\mu\).  Explain why
  \(\SuperD_\alpha S\) is a superfield although
  \(\partial S/\partial\theta_\alpha\) alone is not.

  \SupplementaryExercise{5}{A Superspace Integration by Parts}
    {(Adapted from Stephen P. Martin, \emph{A Supersymmetry Primer},
    version 7, Sections 4.6--4.7)}
    {martin-susy-primer-v7-ch26}
  Prove Eq.~\eqref{eq:26.3.32} by Grassmann degree counting and an ordinary
  spacetime integration by parts.  Deduce Eq.~\eqref{eq:26.3.33} for two
  bosonic superfields.

  \SupplementaryExercise{6}{Solving the Chiral Constraint}
    {(Adapted from David Tong, \emph{Supersymmetric Field Theory},
    Section 3.1)}
    {tong-supersymmetric-field-theory-2022-ch26}
  Show directly that
  \(x_+^\mu=x^\mu+\bar\theta_R\gamma^\mu\theta_L\) and \(\theta_L\)
  are annihilated by \(\SuperD_R\).  Derive the expansion
  \eqref{eq:26.3.21} of a left-chiral scalar superfield and count its
  independent off-shell components.

  \SupplementaryExercise{7}{Multiplying Chiral Superfields}
    {(Adapted from David Tong, \emph{Supersymmetric Field Theory},
    Sections 3.1--3.2)}
    {tong-supersymmetric-field-theory-2022-ch26}
  For two left-chiral multiplets
  \((\phi_i,\psi_{iL},\mathcal F_i)\), find the three components of
  \(\Phi_1\Phi_2\).  Check the symmetry of the fermion bilinear under
  \(1\leftrightarrow2\).

  \SupplementaryExercise{8}{Why an \(\mathcal F\)-Term Is Invariant}
    {(Adapted from Adel Bilal, \emph{Introduction to Supersymmetry},
    Sections 4.2--4.3)}
    {bilal-introduction-susy-2001-ch26}
  Use the component transformations
  \eqref{eq:26.3.15}--\eqref{eq:26.3.17} to show that
  \(\delta[\Phi]_{\mathcal F}\) is a spacetime derivative.  Repeat the
  argument for the \(D\)-component of a general scalar superfield.

  \SupplementaryExercise{9}{Kahler Transformations}
    {(Adapted from Philip C. Argyres, \emph{An Introduction to Global
    Supersymmetry}, Section 1.5.4)}
    {argyres-global-susy-2001-ch26}
  Prove that the action is unchanged by
  \[
    K(\Phi,\Phi^*)\longrightarrow K(\Phi,\Phi^*)+
      h(\Phi)+h(\Phi)^*
  \]
  for any holomorphic \(h\).  Show that the Kahler metric
  \(K_{n\bar m}\) is unchanged.

  \SupplementaryExercise{10}{Renormalizability by Dimension Counting}
    {(Adapted from Stephen P. Martin, \emph{A Supersymmetry Primer},
    version 7, Sections 3.2 and 4.6--4.7)}
    {martin-susy-primer-v7-ch26}
  With \(d(\Phi)=1\), \(d(\SuperD)=1/2\), derive the bounds
  \(d(f)\le3\) and \(d(K)\le2\) for a renormalizable theory.  After removing
  chiral \(D\)-terms and derivative redundancies, show that \(f\) is at
  most cubic and \(K\) is a Hermitian quadratic form.

  \SupplementaryExercise{11}{Eliminating the Chiral Auxiliary Field}
    {(Adapted from Stephen P. Martin, \emph{A Supersymmetry Primer},
    version 7, Sections 3.1--3.2)}
    {martin-susy-primer-v7-ch26}
  For canonical \(K=\sum_n\Phi_n^*\Phi_n\), isolate all terms containing
  \(\mathcal F_n\) in the component Lagrangian.  Complete the square,
  eliminate \(\mathcal F_n\), and obtain the scalar potential.

  \SupplementaryExercise{12}{When Does an \(F\)-Term Vacuum Exist?}
    {(Adapted from Adel Bilal, \emph{Introduction to Supersymmetry},
    Sections 6.1 and 6.3.1)}
    {bilal-introduction-susy-2001-ch26}
  For canonical kinetic terms prove that a classical supersymmetric vacuum
  exists precisely when the equations \(f_n(\phi)=0\) have a common
  solution.  Explain how positivity of \(V=\sum_n|f_n|^2\) enters both
  directions.

  \SupplementaryExercise{13}{Symmetry of the Fermion Mass Matrix}
    {(Adapted from Philip C. Argyres, \emph{An Introduction to Global
    Supersymmetry}, Sections 1.5.5 and 1.6.1)}
    {argyres-global-susy-2001-ch26}
  Expand a holomorphic superpotential to quadratic order around a constant
  background.  Show that the chiral-fermion mass matrix is \(f_{nm}\) and
  is symmetric.  Relate its complex conjugate to the right-chiral mass
  term in four-component notation.

  \SupplementaryExercise{14}{A Moduli-Space Tangent Vector}
    {(Adapted from Philip C. Argyres, \emph{An Introduction to Global
    Supersymmetry}, Sections 1.6.3--1.6.4)}
    {argyres-global-susy-2001-ch26}
  Suppose \(f_n(\phi(s))=0\) along a smooth one-complex-parameter family of
  supersymmetric vacua.  Differentiate with respect to \(s\) and show that
  the tangent vector is a zero mode of \(f_{nm}\).  Identify the associated
  massless chiral multiplet.

  \SupplementaryExercise{15}{The Goldstino Null Eigenvector}
    {(Adapted from Adel Bilal, \emph{Introduction to Supersymmetry},
    Sections 6.1--6.2)}
    {bilal-introduction-susy-2001-ch26}
  At a stationary point of \(V=\sum_n|f_n|^2\), prove
  \(f_{mn}f_n^*=0\).  Conclude that when some \(f_n\ne0\), the normalized
  combination \(f_n^*\psi_{nL}/(\sum|f|^2)^{1/2}\) is massless.

  \SupplementaryExercise{16}{A Minimal O'Raifeartaigh Model}
    {(Adapted from David Tong, \emph{Supersymmetric Field Theory},
    Section 3.4)}
    {tong-supersymmetric-field-theory-2022-ch26}
  Take
  \(f=X(\lambda\Phi^2-\mu^2)+mY\Phi\), with nonzero
  \(m,\lambda,\mu\).  Show that the \(F\)-equations cannot all vanish.
  Find the tree-level vacuum at \(\Phi=Y=0\) in the parameter range where
  it is locally stable, identify the pseudomodulus, and find the goldstino.

  \SupplementaryExercise{17}{The Current as a Moment Map}
    {(Adapted from Adel Bilal, \emph{Introduction to Supersymmetry},
    Section 7.1)}
    {bilal-introduction-susy-2001-ch26}
  Let \(\delta\Phi_n=i\epsilon\mathcal T_{nm}\Phi_m\) be an isometry of
  \(K\).  Show that
  \(\mathcal J=K_n\mathcal T_{nm}\Phi_m\) is real, using only the
  infinitesimal invariance of \(K\).  Explain why it is defined only up to a
  field-independent real constant.

  \SupplementaryExercise{18}{The Canonical \(U(1)\) Current Multiplet}
    {(Adapted from Stephen P. Martin, \emph{A Supersymmetry Primer},
    version 7, Sections 4.4 and 4.7)}
    {martin-susy-primer-v7-ch26}
  For one chiral field with \(K=\Phi^*\Phi\) and
  \(\Phi\mapsto e^{i\epsilon}\Phi\), expand
  \(\mathcal J=\Phi^*\Phi\) through its vector component.  Verify directly
  that this vector is the scalar-plus-fermion Noether current.

  \SupplementaryExercise{19}{Components of a Linear Superfield}
    {(Adapted from Adel Bilal, \emph{Introduction to Supersymmetry},
    Sections 4.1--4.3)}
    {bilal-introduction-susy-2001-ch26}
  Apply \(\SuperD_R^2\) and \(\SuperD_L^2\) to a general real superfield.
  Derive
  \(M=N=0,\ \partial_\mu V^\mu=0,\ 
  \lambda=-\sl{\partial}\omega,\ D=-\Box C\).
  Count the remaining independent off-shell components.

  \SupplementaryExercise{20}{The Superfield Euler--Lagrange Equation}
    {(Adapted from Philip C. Argyres, \emph{An Introduction to Global
    Supersymmetry}, Sections 1.5.4--1.5.5)}
    {argyres-global-susy-2001-ch26}
  Vary the action \eqref{eq:26.6.9} while preserving chirality by writing
  \(\delta\Phi_n=\SuperD_R^2\delta S_n\).  Integrate by parts in superspace
  and derive Eq.~\eqref{eq:26.6.15}, including its factor of \(-4\).

  \SupplementaryExercise{21}{The Kahler Metric in Components}
    {(Adapted from Adel Bilal, \emph{Introduction to Supersymmetry},
    Section 7.1)}
    {bilal-introduction-susy-2001-ch26}
  From the \(D\)-term of a general derivative-free \(K(\Phi,\Phi^*)\),
  show that the scalar kinetic term is
  \(-K_{n\bar m}\partial_\mu\phi_n\partial^\mu\phi_m^*\).
  Prove that positivity requires \(K_{n\bar m}\) to be positive definite.

  \SupplementaryExercise{22}{Auxiliaries with Noncanonical \(K\)}
    {(Adapted from David Tong, \emph{Supersymmetric Field Theory},
    Section 3.3)}
    {tong-supersymmetric-field-theory-2022-ch26}
  Retain the two-fermion terms and derive the algebraic field equation for
  \(\mathcal F_n\) from Eq.~\eqref{eq:26.8.6}.  In the purely bosonic
  limit show that
  \(V=f_n(K^{-1})^{n\bar m}f_{\bar m}^*\).

  \SupplementaryExercise{23}{Kahler Curvature and Four Fermions}
    {(Adapted from Adel Bilal, \emph{Introduction to Supersymmetry},
    Section 7.1)}
    {bilal-introduction-susy-2001-ch26}
  Eliminate the auxiliary fields in a nonlinear sigma model and combine
  the resulting four-fermion term with the last term of
  Eq.~\eqref{eq:26.8.6}.  Show that its coefficient is the Kahler curvature
  tensor \(R_{n\bar m p\bar q}\), up to the spinor-ordering convention.

  \SupplementaryExercise{24}{Conservation of the Supersymmetry Current}
    {(Adapted from Stephen P. Martin, \emph{A Supersymmetry Primer},
    version 7, Sections 3.1--3.2)}
    {martin-susy-primer-v7-ch26}
  Take the divergence of Eq.~\eqref{eq:26.7.10}.  Use the scalar and
  fermion field equations to prove \(\partial_\mu S^\mu=0\), keeping
  \(f_{nm}\) explicit until the final cancellation.

  \SupplementaryExercise{25}{An Identically Conserved Improvement}
    {(Adapted from David Tong, \emph{Supersymmetric Field Theory},
    Sections 3.2--3.3)}
    {tong-supersymmetric-field-theory-2022-ch26}
  Prove without field equations that
  \[
    \partial_\mu\partial_\nu\!
      \left([\gamma^\mu,\gamma^\nu]\chi\right)=0
  \]
  for any spinor \(\chi\).  Explain why the improvement
  \eqref{eq:26.7.18} changes neither the conserved charge nor its
  supersymmetry transformations.

  \SupplementaryExercise{26}{Gamma Trace and Scale Invariance}
    {(Adapted from Philip C. Argyres, \emph{An Introduction to Global
    Supersymmetry}, Sections 1.6.1 and 3.4)}
    {argyres-global-susy-2001-ch26}
  Use Euler's theorem for a homogeneous cubic superpotential to show that
  the right-hand side of Eq.~\eqref{eq:26.7.19} vanishes.  Check the
  converse for a polynomial superpotential with canonical dimensions:
  identify which mass and linear terms obstruct gamma-tracelessness.

  \SupplementaryExercise{27}{Inverting the Supercurrent Relation}
    {(Adapted from Adel Bilal, \emph{Introduction to Supersymmetry},
    Sections 4.3 and 7.1)}
    {bilal-introduction-susy-2001-ch26}
  Starting from
  \(S^\mu=-2\omega^\mu+2\gamma^\mu\gamma^\nu\omega_\nu\),
  solve algebraically for \(\omega_\mu\) in terms of \(S_\mu\).  Use only
  \(\gamma_\mu\gamma^\mu=4\).

  \SupplementaryExercise{28}{Parity Exchanges Chiralities}
    {(Adapted from Philip C. Argyres, \emph{An Introduction to Global
    Supersymmetry}, Sections 1.3.2 and 1.5.3)}
    {argyres-global-susy-2001-ch26}
  Apply \(x\mapsto\Lambda_Px,\ \theta\mapsto-i\beta\theta\) to
  \(x_\pm\) and to the chiral projectors.  Prove that parity maps a
  left-chiral scalar superfield into a right-chiral one and write the
  induced map on \((\phi,\psi_L,\mathcal F)\).

  \SupplementaryExercise{29}{Time Reversal Preserves Chirality}
    {(Adapted from Philip C. Argyres, \emph{An Introduction to Global
    Supersymmetry}, Sections 1.3.2 and 1.5.3)}
    {argyres-global-susy-2001-ch26}
  Let \(B_T=-\gamma_5\mathcal C\) be the Majorana time-reversal matrix.
  Taking account of antiunitarity, show that the time reverse of
  \(\SuperD_R\Phi=0\) again obeys a right-superderivative constraint rather
  than a left-superderivative constraint.  Explain the contrast with
  parity in terms of massless helicity.

  \SupplementaryExercise{30}{A Dimension-Five Operator Basis}
    {(Adapted from Stephen P. Martin, \emph{A Supersymmetry Primer},
    version 7, Sections 4.6--4.7 and 4.10)}
    {martin-susy-primer-v7-ch26}
  For one neutral left-chiral scalar superfield, classify all Lorentz-scalar
  supersymmetric operators of component dimension five.  Work modulo
  complex conjugation, ordinary and superspace integration by parts, and
  chiral \(D\)-terms.  Include the higher-derivative class.
\end{enumerate}
